% Architecture hexagonale de la plateforme (figure 3.1)
\begin{tikzpicture}[x=1cm, y=1cm]

  % ── Hexagone extérieur : la couronne des ports ────────────
  \node[regular polygon, regular polygon sides=6, draw=black, fill=figGrisTresClair,
        minimum size=76mm, inner sep=0pt, rotate=0] (hexext) at (0,0) {};

  % ── Hexagone intérieur : le domaine métier ────────────────
  \node[regular polygon, regular polygon sides=6, draw=black, fill=white,
        minimum size=50mm, inner sep=0pt] (hexint) at (0,0) {};

  \node[font=\sffamily\scriptsize\bfseries, align=center] at (0,0.55) {Domaine métier};
  \node[font=\sffamily\tiny, align=center] at (0,-0.05)
       {Entités\\Objets valeur\\Règles de gestion};
  \node[font=\sffamily\tiny\itshape, align=center] at (0,-1.02)
       {ne dépend de rien};

  % Légende de la couronne, posée à l'extérieur avec un renvoi
  \node[font=\sffamily\tiny\itshape, align=right, anchor=south east]
       (portlbl) at (-3.05, 3.55) {Couronne des ports :\\interfaces déclarées\\par le domaine};
  \draw[dependance] (portlbl.south east) -- (-1.75,2.45);

  % ── Ports posés sur la couronne ───────────────────────────
  \foreach \a/\lbl in {30/{Décision}, 90/{Politique}, 150/{Exécution},
                       210/{Connaissance}, 270/{Audit}, 330/{Intégrations}}
    {
      \node[draw=black, fill=white, rectangle, inner sep=1.6pt,
            font=\sffamily\tiny] at ({\a}:3.15) {\lbl};
    }

  % ── Adaptateurs entrants ──────────────────────────────────
  \node[fbox, minimum width=30mm, align=center] (ae1) at (-6.6, 1.55)
       {Couche temps réel\\(entrée conversationnelle)};
  \node[fbox, minimum width=30mm, align=center] (ae2) at (-6.6,-0.15)
       {Interface applicative\\métier};
  \node[fbox, minimum width=30mm, align=center] (ae3) at (-6.6,-1.85)
       {Interfaces web};

  \node[etiquettezone, anchor=south] at (-6.6, 2.45) {Adaptateurs entrants};
  \draw[zone] ($(ae1.north west)+(-0.18,0.30)$) rectangle ($(ae3.south east)+(0.18,-0.30)$);

  % ── Adaptateurs sortants ──────────────────────────────────
  \node[fbox, minimum width=30mm, align=center] (as1) at (6.6, 2.15) {Persistance};
  \node[fbox, minimum width=30mm, align=center] (as2) at (6.6, 0.95) {Index documentaire};
  \node[fbox, minimum width=30mm, align=center] (as3) at (6.6,-0.25) {Gestion de tickets};
  \node[fbox, minimum width=30mm, align=center] (as4) at (6.6,-1.45) {Systèmes de l'opérateur};
  \node[fbox, minimum width=30mm, align=center] (as5) at (6.6,-2.65) {Journal d'audit};

  \node[etiquettezone, anchor=south] at (6.6, 2.95) {Adaptateurs sortants};
  \draw[zone] ($(as1.north west)+(-0.18,0.30)$) rectangle ($(as5.south east)+(0.18,-0.30)$);

  % ── Dépendances : toutes dirigées vers le centre ──────────
  \foreach \n in {ae1,ae2,ae3}{ \draw[flux] (\n.east) -- ($(hexext.west)+(-0.05,0)$); }
  \foreach \n in {as1,as2,as3,as4,as5}{ \draw[flux] (\n.west) -- ($(hexext.east)+(0.05,0)$); }

  \node[note, text width=88mm, anchor=north, align=center] at (0,-4.30)
       {Le sens des dépendances est inversé : l'infrastructure se conforme au contrat défini par le\\domaine. Remplacer un fournisseur revient à écrire un adaptateur, sans toucher aux règles.};

\end{tikzpicture}
